\chapter{Revisão de Literatura, Inovação e Lacunas}

A literatura é organizada em três linhas relativamente maduras, porém pouco integradas: antenas e arrays conformais de banda larga; radomes funcionais e superfícies seletivas em frequência; e detecção passiva, localização e identificação de emissores. Estudos de radiogoniometria conformal mostram que acoplamento mútuo, estrutura de suporte e distorção do radome pertencem ao manifold calibrado do array, não sendo apenas questões mecânicas \cite{caratelli2011,liberal2011}.

Na literatura de engenharia, \emph{radome} é a contração de \emph{radar dome} e designa a estrutura que protege uma antena contra o ambiente, sob o compromisso entre integridade mecânica e transparência eletromagnética \cite{qamar2020}. O artigo disponibilizado em \texttt{bibliography/} avalia esse compromisso em sistemas ativos de radar e comunicação e mostra que perda de transmissão, incidência oblíqua, descasamento copolar e despolarização precisam ser medidos com a antena e o radome integrados. A definição estrutural não obriga, contudo, que a antena protegida transmita. A arquitetura aqui proposta especializa o conceito como um domo de recepção: não há iluminador RF próprio no nó, e todas as cadeias associadas à vigilância são receptoras.

A pesquisa de radomes funcionais demonstra transmissão em banda combinada com absorção ou difusão fora das bandas selecionadas \cite{costa2012,lv2021,omar2018,sheng2024}. Isso cria uma tensão de projeto no monitoramento aberto do espectro: um radome otimizado para baixa assinatura ou transmissão estreita pode rejeitar sinais que o receptor deveria observar.

A pesquisa de localização passiva utiliza AOA, TDOA e FDOA, com precisão governada por geometria, sincronização, propagação e carga computacional \cite{liu2019,pine2021}. Fingerprinting de RF e detecção de anomalias em conjunto aberto podem auxiliar a caracterização de emissores, mas classificadores fechados e condições limitadas de treinamento não devem ser tratados como soluções universais \cite{soltanieh2020,jagannath2022,sun2022}.

Essa separação permite duas classes de inferência passiva que não devem ser confundidas. Sinais diretos de um emissor podem ser localizados por triangulação de AOA ou por multilateração com TDOA/FDOA. Para alvos que não emitem, a rede usa iluminadores de oportunidade externos e compara o canal direto de referência com os canais de vigilância que contêm reflexões biestáticas ou multiestáticas. Em ambos os casos os radomes apenas recebem; a posição resulta da geometria e da fusão calibrada entre nós, não de uma transmissão de interrogação produzida pelo sistema.

A lacuna do projeto é a integração ponta a ponta de uma abertura conformal multifaixa, recepção passiva distribuída, calibração polarimétrica e temporal, processamento hierárquico de eventos e validação realista em sítios montanhosos. A inovação é sistêmica e deve ser demonstrada por medições reproduzíveis.

\section{Inovação arquitetural candidata: abertura externa e célula absorvedora}
A combinação proposta de Yagis VHF/UHF expostas além do casco com células piramidais internas individualmente blindadas e revestidas por absorvedor constitui uma lacuna de integração a investigar. A antena externa evita que o caminho útil atravesse a cobertura dielétrica; atrás de cada face, a célula busca confinar a eletrônica sensível, limitar acoplamento entre canais e dissipar a parcela de energia que penetra na abertura sem ser entregue ao receptor. Em princípio, isso pode reunir recepção eficiente de iluminadores existentes, menor auto-interferência e menor reradiação pela cavidade. A literatura de radomes absorventes e difusivos demonstra que transmissão seletiva, absorção e controle de espalhamento podem ser co-projetados, mas não valida esta geometria específica \cite{costa2012,lv2021}.

Essa proposta é registrada como inovação arquitetural a comprovar, não como invisibilidade demonstrada. Uma parede Faraday metálica sem material dissipativo reflete a onda; além disso, Yagis, boom, bordas, juntas e cabos externos continuam espalhando energia. A hipótese de baixa observabilidade exige especificar material e espessura do absorvedor e medir, em função de frequência, polarização e ângulo, o balanço entre potência recebida, refletida, absorvida e reradiada, além de $S_{11}$ da abertura, eficácia de blindagem, acoplamento entre células e RCS monoestática e biestática com e sem revestimento. O caráter passivo do sistema reduz sua assinatura de emissão porque não há transmissor próprio, mas não implica invisibilidade a sensores externos.

\section{Cobertura das lacunas pela arquitetura atual}
O projeto atual resolve várias lacunas no nível arquitetural. O manifold calibrado inclui o radome, a geometria da face, o acoplamento mútuo e a orientação local das antenas; a abertura híbrida atribui uma Yagi VHF externa de polarização única e uma Yagi UHF externa menor de polarização única a módulos RF independentes; módulos dual-porta na mesma faixa são obrigatórios onde houver síntese polarimétrica; cada antena possui cadeia ADC/ASIC blindada; o FFASIC executa a fusão da face; e o projeto térmico e estrutural define placas frias, sensores, caminhos internos de carga e gates ambientais. Essas são soluções de projeto, não afirmações de desempenho já medido. O trabalho restante é quantificar perda de inserção, fase e atraso de grupo, margens estruturais, gradientes térmicos, isolamento EMC e covariância de localização em ensaios representativos.

\section{Cobertura da lacuna pela abertura separada por faixas}
A revisão identificou a ausência de um tratamento ponta a ponta da integração da abertura, polarização, calibração e efeitos realistas da plataforma. A arquitetura atual aborda isso no nível de projeto ao separar integração mecânica entre faixas de polarimetria dentro da mesma faixa. As Yagis VHF e UHF fornecem observações espectrais independentes; somente um módulo dual-porta dentro de uma faixa pode produzir grandezas calibradas de Jones, Stokes, RHCP ou LHCP. A questão restante é quantitativa: ganho, largura de banda, isolamento, polarização cruzada dos módulos dual-porta válidos, efeito da plataforma e estabilidade devem confirmar o projeto.
